% =============================================================================
\section{Tarefa 2 --- Trajetória de swing}
\label{sec:traj}
% =============================================================================

\subsection{Analogia com o Laboratório 1: o swing desenhado nas juntas}
\label{sec:perfis}

No Lab~1 o swing era um pêndulo duplo~\cite{cochran1968} movido por três
juntas: o tronco roda, o braço sobe e desce em torno do ombro e o pulso fica
armado no topo e só solta perto do impacto. Aqui o robô faz de jogador e a
pinça faz de taco, com a ponta a bater na bola. A analogia é feita junta a
junta: a base (junta~1) é o \textbf{tronco}, a junta~2 é o \textbf{ombro} e a
junta~5 é o \textbf{pulso}. O cotovelo fica fixo e as juntas 4 e 6 a zero, de
modo que o braço se move sempre no plano vertical que a base roda.

O movimento descreve-se com os dois ângulos do Lab~1, $\varphi(t)$ (braço,
$\varphi=0$ na bola) e $\psi(t)$ (pulso), com as mesmas três fases e leis de
movimento (Tabela~\ref{tab:fases}); $s(\tau) = 10\tau^3 - 15\tau^4 + 6\tau^5$
é o polinómio de grau~5 com velocidade e aceleração nulas nos extremos.

\begin{table}[H]
\centering
\caption{Fases do swing. $\varphi_{top} = \ang{-90}$,
$\varphi_{end} = \ang{50}$, $\psi_{top} = \ang{30}$; $u$ é o tempo contado
desde o início de cada fase.}
\label{tab:fases}
\small
\begin{tabularx}{\textwidth}{@{}lcLL@{}}
\toprule
\rowcolor{tabcab}
Fase & Duração & Braço $\varphi$ & Pulso $\psi$ \\
\midrule
Backswing & $T_b = \SI{4.5}{\second}$ &
  $\varphi_{top}\,s(u/T_b)$: parte e chega parado &
  $\psi_{top}\,s(u/T_b)$: arma \\
Downswing & $T_d = \SI{4.0}{\second}$ &
  $\varphi_{top}\bigl(1-(u/T_d)^2\bigr)$: aceleração constante, velocidade máxima na bola &
  armado até $0{,}6\,T_d$; depois solta com $s$ e chega a 0 no impacto \\
Follow-through & $T_f = \SI{2.22}{\second}$ &
  $\omega_{imp}\,u - \tfrac12 (\omega_{imp}/T_f)\,u^2$: desaceleração constante até $\varphi_{end}$ &
  0 \\
\bottomrule
\end{tabularx}
\end{table}

As juntas seguem estes ângulos (em graus):
\begin{equation}
q_1 = -0{,}8\,\varphi,\qquad
q_2 = 48{,}8 - 35\,(1-\cos\varphi),\qquad
q_3 = 80,\qquad
q_5 = 29{,}5 + \psi,\qquad
q_4 = q_6 = 0 .
\label{eq:juntas}
\end{equation}
A posição e a orientação da pinça em função do tempo, que é o que o enunciado
pede, obtêm-se compondo~\eqref{eq:juntas} com a cinemática direta da Tarefa~1:
\begin{equation}
\begin{bmatrix} R(t) & p(t)\\ 0 & 1\end{bmatrix} = \T{0}{6}\bigl(q(t)\bigr)\,T_z(L).
\label{eq:pose}
\end{equation}
No endereço ($\varphi = \psi = 0$) a ponta da pinça está em
$b = (347;\,0;\,80)$~mm, que é a posição da bola, com a pinça vertical (a
menos de \ang{2}). No topo do backswing está em $(95;\,293;\,281)$~mm e no fim
em $(287;\,-240;\,159)$~mm (Figuras~\ref{fig:cart} e~\ref{fig:3d}). No
impacto, $\omega_{imp} = -2\varphi_{top}/T_d = \SI{0.785}{\radian\per\second}$
e a ponta passa na bola a \SI{218}{\milli\meter\per\second}, na horizontal e
ao longo de $-\hat y_0$: bate-lhe de lado, como a face de um taco.

\begin{figure}[H]
\centering
\begin{minipage}[b]{0.55\textwidth}
\centering
\fig{Photos/traj_golfista_perto_cartesiana.png}{\textwidth}
\caption{Posição, orientação e velocidade da ponta da pinça e do flange, com
o limite do fabricante.}
\label{fig:cart}
\end{minipage}\hfill
\begin{minipage}[b]{0.43\textwidth}
\centering
\fig{Photos/traj_golfista_perto_3d.png}{\textwidth}
\caption{Caminho da ponta e do flange; a cinzento, a pinça em instantes
sucessivos; os círculos marcam o cilindro da base.}
\label{fig:3d}
\end{minipage}
\end{figure}

Quatro escolhas justificam-se:

\begin{itemize}[itemsep=3pt]
  \item \textbf{Desenho nas juntas e não no espaço cartesiano.} Uma primeira
        versão impunha à ponta um arco cartesiano e deixava a inversa escolher
        as juntas; o antebraço tinha de rodar \ang{121} e o robô não seguia o
        perfil à velocidade planeada (Secção~\ref{sec:v1}). Aqui cada junta
        tem um papel e um perfil suave, e só três se movem.
  \item \textbf{Velocidade máxima na bola.} O backswing é mais lento do que o
        downswing, e o máximo de velocidade de toda a trajetória ocorre no
        impacto, como num swing real.
  \item \textbf{Pulso longe da singularidade.} Com as juntas 4 e 6 a zero, o
        pulso é singular em $q_5 = 0$. O pulso arma no sentido que afasta
        $q_5$ de zero: fica entre \ang{29.5} e \ang{59.5}.
  \item \textbf{Cotovelo dobrado a \ang{80}.} O braço usa \SI{77}{\percent} do
        comprimento esticado, longe da singularidade do cotovelo.
\end{itemize}

\subsection{Viabilidade}
\label{sec:viab}

As velocidades e acelerações calculam-se por diferenças centrais sobre os 1073
pontos da trajetória (\SI{10.72}{\second} a \SI{100}{\hertz}). A
Tabela~\ref{tab:viab} resume as três condições do enunciado, com
\SI{20}{\percent} de margem sobre os limites do
fabricante~\cite{ufactory_lite6,ufactory_specs}. O alcance mede-se do ombro
ao centro do pulso, que não pode passar o braço esticado,
$200 + \sqrt{87^2 + 227{,}6^2} = \SI{443.7}{\milli\meter}$; perto desse valor
o cotovelo entra em singularidade, por isso só se aceita até
\SI{95}{\percent}.

\begin{table}[H]
\centering
\caption{Verificação de viabilidade (\texttt{check\_feasibility} em
\texttt{src/trajectory.py}).}
\label{tab:viab}
\small
\begin{tabularx}{\textwidth}{@{}lLccc@{}}
\toprule
\rowcolor{tabcab}
Condição & O que se mede & Limite & Aceite até & Trajetória \\
\midrule
\rowcolor{tabgrupo}\multicolumn{5}{@{}l}{\emph{(i) Velocidade e aceleração}} \\
Ponta da pinça & $\max\lVert\dot p\rVert$ & \SI{500}{\milli\meter\per\second} & 400 & \textbf{218} \\
Flange         & $\max\lVert\dot p_f\rVert$ & \SI{500}{\milli\meter\per\second} & 400 & \textbf{216} \\
Juntas         & $\max_i\lvert\dot q_i\rvert$ & \SI{180}{\degree\per\second} & 144 & \textbf{36} (junta 1) \\
Acel.\ das juntas & $\max_i\lvert\ddot q_i\rvert$ & \SI{1145}{\degree\per\second\squared} & 916 & \textbf{68} (junta 5) \\
\midrule
\rowcolor{tabgrupo}\multicolumn{5}{@{}l}{\emph{(ii) Alcance}} \\
Ombro ao centro do pulso & $\max\lVert p_w - p_{\{2\}}\rVert$ & \SI{443.7}{\milli\meter} & 421,5 & \textbf{343} \\
\midrule
\rowcolor{tabgrupo}\multicolumn{5}{@{}l}{\emph{(iii) Limites das juntas}} \\
Folga ao limite mais próximo & $\min_i$ (Tabela~\ref{tab:juntas}) & --- & $\geq\ang{5}$ & \textbf{\ang{64}} (junta 5) \\
\bottomrule
\end{tabularx}
\end{table}

\begin{figure}[H]
\centering
\begin{minipage}[c]{0.42\textwidth}
\centering
\small
\captionof{table}{Limites das juntas~\cite{ufactory_specs} e gama usada.}
\label{tab:juntas}
\begin{tabular}{@{}ccc@{}}
\toprule
\rowcolor{tabcab}
Junta & Limites & Gama usada \\
\midrule
1 & $\pm\ang{360}$         & \ang{-40} a \ang{72}  \\
2 & $\pm\ang{150}$         & \ang{14} a \ang{49}  \\
3 & \ang{-3.5} a \ang{300} & \ang{80} (fixa)  \\
4 & $\pm\ang{360}$         & \ang{0} (fixa) \\
5 & $\pm\ang{124}$         & \ang{29.5} a \ang{59.5}  \\
6 & $\pm\ang{360}$         & \ang{0} (fixa) \\
\bottomrule
\end{tabular}
\end{minipage}\hfill
\begin{minipage}[c]{0.55\textwidth}
\centering
\fig{Photos/traj_golfista_perto_juntas.png}{\textwidth}
\captionof{figure}{Ângulos e velocidades das juntas; a tracejado, o limite
de \SI{180}{\degree\per\second}.}
\label{fig:juntas}
\end{minipage}
\end{figure}

\subsection{Ausência de contacto perigoso: simulação no MuJoCo}
\label{sec:contacto}

A verificação anterior garante que a ponta, o flange e as origens de todos os
referenciais ficam a mais de \SI{30}{\milli\meter} da mesa (mínimo:
\SI{80}{\milli\meter}, na bola) e a mais de \SI{80}{\milli\meter} do eixo da
base (mínimo: \SI{298}{\milli\meter}). Como isto só olha para pontos, antes de
ir ao robô reproduziu-se a trajetória das juntas no MuJoCo~\cite{todorov2012},
com o modelo oficial do Lite~6 com pinça~\cite{menagerie2022}, e mediu-se em
cada instante a distância entre as malhas de colisão de todos os pares de
elos, incluindo o corpo e os dedos da pinça, e entre o robô e a mesa
(\texttt{src/mujoco\_preview.py}). As malhas são invólucros convexos, logo as
distâncias são conservadoras.

\begin{figure}[H]
\centering
\fig{Photos/mjgp_sequencia_face_on.png}{\textwidth}
\caption{O swing no modelo 3D do robô. A vermelho, o caminho da ponta da
pinça; a esfera é a bola.}
\label{fig:seq}
\end{figure}

\begin{figure}[H]
\centering
\fig{Photos/mjgp_folgas.png}{0.6\textwidth}
\caption{Folga mínima ao longo do swing: robô (com pinça) à mesa e entre elos
não adjacentes. As linhas verticais marcam o topo do backswing e o impacto.}
\label{fig:folgas}
\end{figure}

O ponto mais baixo do robô é a pinça, a \SI{82}{\milli\meter} da mesa no
endereço e no impacto. Entre elos, a menor distância é de
\SI{25}{\milli\meter} (corpo da pinça e elo~4) e é praticamente a distância
construtiva entre os dois; o par seguinte (elos 2 e 4) nunca fica a menos de
\SI{35}{\milli\meter}. O modelo 3D, construído pelo fabricante a partir do
desenho do robô, e a nossa cinemática direta coincidem a menos de
\SI{0.7}{\milli\meter} no flange e \SI{1.9}{\milli\meter} na ponta da pinça ao
longo de toda a trajetória, o que é uma verificação independente da Tarefa~1.
A simulação serviu também para escolher a trajetória: várias versões anteriores
foram postas de parte depois de vistas no modelo, sem chegarem ao robô.
